% 张三慧第4版 V2第27章 11题独立LaTeX解析
% V2 PDF386--387，S PDF365--369；候选稿，需原图及专家复核。
\section*{第27章\quad 薛定谔方程：独立解析}
\paragraph{习题 27.1}
无限深势阱的谱为
\[
E_n=\frac{\pi^2\hbar^2 n^2}{2ma^2}
=\frac{h^2n^2}{8ma^2},\quad n=1,2,\ldots.
\]
\(m=10^{-14}\,\mathrm g=10^{-17}\,\mathrm{kg}\)，
\(a=10^{-5}\,\mathrm m\)。于是
\[
\boxed{E_{100}\simeq5.488\times10^{-37}\,\mathrm J},
\quad
\boxed{E_{101}\simeq5.598\times10^{-37}\,\mathrm J},
\]
\[
\boxed{\Delta E
=\frac{h^2}{8ma^2}(101^2-100^2)
\simeq1.103\times10^{-38}\,\mathrm J}.
\]
极微小能级间隔说明在细胞尺度、较大质量下
能谱对实验几乎连续。

\paragraph{习题 27.2}
取氧气分子 \(m=32.0\times10^{-3}/N_A
\simeq5.314\times10^{-26}\,\mathrm{kg}\)，
\(a=.10\,\mathrm m\)：
\[
\boxed{E_1=\frac{h^2}{8ma^2}
\simeq1.033\times10^{-40}\,\mathrm J}.
\]
令 \(E_n=n^2E_1=(3/2)k_BT\)，则
\[
\boxed{n\simeq\sqrt{\frac{3k_BT}{2E_1}}
\simeq7.76\times10^9}.
\]
相邻能级差
\[
\boxed{\Delta E
=E_1[(n+1)^2-n^2]=E_1(2n+1)
\simeq1.60\times10^{-30}\,\mathrm J}.
\]
在如此巨大的 \(n\) 下，相邻能级不可实验分辨。
配套S PDF366将该结果写为约
\(1.6\times10^{-30}\,\mathrm J\)，符合精算。

\paragraph{习题 27.3}
本题是\emph{斜底}无限深势阱，
不能用普通平底井的正弦函数直接绘图。
对总能 \(E_5\)，局部经典动量
\(p(x)=\sqrt{2m[E_5-U(x)]}\)，
只在 \(E_5>U(x)\) 的允许区有振荡解。
近似局部半波长
\(\lambda(x)/2=h/[2p(x)]\)：
势底由左往右抬升时，右侧 \(E_5-U\) 较小，
波长较长，\emph{同一概率流正规化下}
波函数的局部振幅亦较大。
左侧振荡更密、振幅较小；
在垂直壁处 \(\psi=0\)，有限斜坡越过
\(E_5\) 的禁阻区中呈指数衰减。
\(n=5\) 的非简并一维束缚本征态按节点定理
在允许区域有4个内部节点。
原图27.14(a)势阱右侧仍是垂直无穷壁，
波函数须在右端再次降至零；
(b)斜底右侧超出 \(E_5\) 后的经典禁阻区
应是衰减尾部，不是继续周期正弦。
\textit{确切节点坐标与WKB相位修正取决于
图27.14所画斜坡，仍需逐像素重新制图。}

\paragraph{习题 27.4}
归一化基态
\(\psi_1(x)=\sqrt{2/a}\sin(\pi x/a)\)。
因此靠左边壁四分之一井宽内的概率
\[
P=\frac2a\int_0^{a/4}\sin^2\left(\frac{\pi x}a\right)\,dx
=\left[\frac xa-\frac{\sin(2\pi x/a)}{2\pi}
\right]_{0}^{a/4}.
\]
\[
\boxed{P=\frac14-\frac1{2\pi}
\simeq0.0908451}.
\]
靠右边壁同长度区间的概率由对称性相等；
不能按均匀分布误答25\%。

\paragraph{习题 27.5}
\(\psi_n(x)=\sqrt{2/a}\sin(n\pi x/a)\)。
由对称性
\[
\boxed{\langle x\rangle
=\frac2a\int_0^a x\sin^2(n\pi x/a)\,dx
=\frac a2}.
\]
再将 \(\sin^2z=(1-\cos2z)/2\) 代入，
积分得
\[
\boxed{\langle x^2\rangle
=\frac2a\int_0^a x^2\sin^2(n\pi x/a)\,dx
=a^2\left(\frac13-\frac1{2\pi^2n^2}\right)}.
\]
故方差
\(\boxed{(\Delta x)^2
=a^2(1/12-1/(2\pi^2n^2))}\)。

\paragraph{习题 27.6}
两定态
\(\Psi_n(x,t)=\sqrt{2/a}\sin(n\pi x/a)
\exp(-iE_nt/\hbar)\)，
因此
\[
\int_0^a\Psi_m^*\Psi_n\,dx
=\frac2a e^{i(E_m-E_n)t/\hbar}
\int_0^a\sin(m\pi x/a)\sin(n\pi x/a)\,dx.
\]
积化和差后
\[
\int_0^a\sin(m\pi x/a)\sin(n\pi x/a)\,dx
=\frac a2\,\delta_{mn},
\]
故当 \(m\ne n\)
\[
\boxed{\langle\Psi_m,\Psi_n\rangle=0}.
\]
这种正交性对一般自伴哈密顿算符的
不同本征值对应的本征态成立。

\paragraph{习题 27.7}
势阱宽度为外加控制参数 \(a\) 时，
定态能量
\(E_n=\pi^2\hbar^2n^2/(2ma^2)\)。
粒子对\emph{右端活动壁}的平均向外压力
从虚功得
\[
\boxed{F=-\frac{\partial E_n}{\partial a}
=\frac{\pi^2\hbar^2n^2}{ma^3}>0}.
\]
这不是两端壁作用力的净矢量和：
两壁分别向外受大小相等的力，
粒子在定态下整体平均动量为零。

\paragraph{习题 27.8}
井壁处波函数为零，故
\[
a=n\frac{\lambda_n}{2},\qquad
\lambda_n=\frac{2a}{n},\qquad n=1,2,\ldots.
\]
德布罗意关系 \(p_n=h/\lambda_n=nh/(2a)\)。
非相对论粒子能量 \(E_n=p_n^2/(2m)\)：
\[
\boxed{E_n=\frac{h^2n^2}{8ma^2}
=\frac{\pi^2\hbar^2n^2}{2ma^2}}.
\]
量子数 \(n=0\) 对应恒零波函数，
不可作为归一化态。

\paragraph{习题 27.9}
立方体各壁坐标为 \(0,a\)，
驻波沿三个轴满足
\(p_j=n_jh/(2a)\)。
\[
E=\frac{p_x^2+p_y^2+p_z^2}{2m}
=\frac{h^2}{8ma^2}(n_x^2+n_y^2+n_z^2).
\]
所以
\[
\boxed{E_{n_xn_yn_z}
=\frac{\pi^2\hbar^2}{2ma^2}
(n_x^2+n_y^2+n_z^2)},
\quad n_x,n_y,n_z\ge1.
\]
最小能量为
\(\boxed{E_{111}=3h^2/(8ma^2)}\)，
三维零点能非零；不同排列量子数
可带来几何简并。

\paragraph{习题 27.10}
谐振子势能 \(V(x)=m\omega^2x^2/2\)。
将 \(\psi_0=Ae^{-\alpha x^2}\)
求导：
\(\psi_0''=(4\alpha^2x^2-2\alpha)\psi_0\)。
定态方程
\[
-\frac{\hbar^2}{2m}\psi_0''+
\frac12m\omega^2x^2\psi_0=E_0\psi_0
\]
给出
\[
\left(\frac{\hbar^2\alpha}{m}-E_0\right)
+\left(\frac{m\omega^2}2-
\frac{2\hbar^2\alpha^2}{m}\right)x^2=0.
\]
对任意 \(x\) 恒等成立，且
\(\alpha>0\) 以保证可归一化，
\[
\boxed{\alpha=\frac{m\omega}{2\hbar}},
\qquad
\boxed{E_0=\frac{\hbar\omega}{2}=\frac{h\nu}{2}}.
\]
负的 \(\alpha\) 会使波函数不归一化，不是基态。

\paragraph{习题 27.11}
\(\omega=\sqrt{k/m}\)，
振动频率 \(\nu=\omega/(2\pi)\)；
相邻能级间距
\[
\Delta E=\hbar\omega
=\hbar\sqrt{\frac{1.13\times10^3}{1.67\times10^{-27}}}
\simeq0.5414\,\mathrm{eV}.
\]
本征谱为
\[
\boxed{E_n=(n+\tfrac12)\,0.5414\,\mathrm{eV}},
\qquad n=0,1,2,\ldots .
\]
因此基态 \(E_0\simeq0.2707\,\mathrm{eV}\)，
相邻能级向下跃迁发射光子
\[
\boxed{E_\gamma=\Delta E\simeq0.5414\,\mathrm{eV}},
\quad
\boxed{\lambda=\frac{hc}{\Delta E}
\simeq2.290\times10^{-6}\,\mathrm m
=2290\,\mathrm{nm}}.
\]
这些是题设单一等效质量 \(m\) 的谐振子模型；
若引入真实H$_2$约化质量与振动非谐性，
频率会改变。
